\documentclass[11pt,a4paper]{article}
\usepackage[T1]{fontenc}
\usepackage{lmodern}
\usepackage[margin=1in]{geometry}
\usepackage{amsmath,amssymb,amsthm,mathtools}
\usepackage{booktabs,array,enumitem,microtype,xcolor}
\usepackage{needspace}
\usepackage{hyperref}
\hypersetup{colorlinks=true,linkcolor=blue!45!black,citecolor=blue!45!black,
 urlcolor=blue!45!black,
 pdftitle={Ordered-instance allocation transfer with exact extra-copy preservation},
 pdfauthor={Dev D. Goyal},
 pdfsubject={Finite additive allocation and a Hall matching construction}}
\setlength{\parindent}{0pt}
\setlength{\parskip}{0.35em}
\newtheorem{theorem}{Theorem}
\newtheorem{corollary}[theorem]{Corollary}
\newtheorem{lemma}[theorem]{Lemma}
\newtheorem{proposition}[theorem]{Proposition}
\newcommand{\mult}{\operatorname{mult}}
\newcommand{\extra}{\operatorname{extra}}
\newcommand{\copies}{\operatorname{copies}}
\newcommand{\nbr}{\mathcal N}
\newcommand{\ranks}{[m]}
\title{Ordered-instance allocation transfer\\with exact extra-copy preservation}
\author{Dev D. Goyal}
\date{28 October 2026}
\begin{document}
\maketitle

\paragraph{Autonomous AI research disclosure.}
This work is fully autonomous agentic research. Automated systems produced the research, proofs, formalization, and manuscript preparation. Dev D. Goyal is responsible for publication. No expert validation or peer review is claimed.

\begin{abstract}
An ordered allocation assigns ranks in each agent's decreasing list of good
values. Transferring such an allocation to named goods must preserve utilities
while respecting the allowed copies of each good. For every finite set of
goods, every nonempty finite agent set, and nonnegative additive real
valuations, an arbitrary set-valued ordered allocation admits a transfer
that weakly increases every agent's utility, preserves the total number of
extra copies exactly, respects the same uniform per-good extra-copy cap,
and allocates every good. Unused input ranks are permitted. The proof
completes the input and uses a balanced bipartite matching graph. Incidence
ports prevent repeated goods within an agent's bundle, and primary copy
vertices force coverage. The explicit graph also gives a polynomial-time
transformation. The accompanying Lean formalization covers the existence
and utility statements; executable runtime analysis is separate.
\end{abstract}

\section{The transfer problem}
Sorting each agent's values turns a named-good instance into an ordered
instance. The rank labels are shared, but a particular rank can name a
different real good for each agent. An ordered allocation is therefore
useful only if it can be transferred back without reducing utilities.
When goods can be copied, the transfer must also control both the total
number of extra copies and the number of extra copies of each good.

The supplied source context attributes an ordered-instance reduction question
to Akrami et al., identified by arXiv:2502.09377v3~\cite{Akrami}.
The primary allocation definition, copy convention, restricted theorem,
and current source version were not verified for this note. The definitions
below are therefore explicit local definitions, and no verified strengthening
or resolution of a published question is claimed. If the source conventions
agree, this theorem may be compared with that motivating question.
No first-priority claim is made. Hall's marriage theorem is the standard
matching result used in the proof.

The main difficulty is not just to find enough eligible goods. An agent
must receive each good at most once, each good has a multiplicity cap, and
all goods must be used to preserve the extra-copy count exactly after
completion. The matching graph enforces these conditions together.

\section{Definitions and the main theorem}
Let $N$ be a nonempty finite set of agents, with $n=|N|\ge1$. Let $M$ be a
finite set of goods, with $m=|M|\ge0$, and write
$\ranks=\{1,\ldots,m\}$, with $[0]=\varnothing$.
Agent $i$ has a nonnegative real weight $v_i(g)$ for each $g\in M$.
Its value for a bundle is additive:
\[
 u_i(B_i)=\sum_{g\in B_i}v_i(g).
\]
For every agent, fix a bijection $r_i:M\to\ranks$ such that
\[
 w_i(a):=v_i(r_i^{-1}(a))
 \quad\text{is nonincreasing in }a.
\]
Ties can be ordered arbitrarily. The ordered utility of $A_i\subseteq\ranks$
is $U_i(A_i)=\sum_{a\in A_i}w_i(a)$.

An allocation on a finite ground set $G$ is a family of sets
$X=(X_i)_{i\in N}$ with $X_i\subseteq G$. Different agents may receive the
same good. Each individual bundle is a set and contains no repeated good.
Define
\begin{align*}
 \mult_X(g)&=|\{i\in N:g\in X_i\}|,\\
 \extra_X(g)&=\max\{\mult_X(g)-1,0\},\\
 \copies(X)&=\sum_{g\in G}\extra_X(g).
\end{align*}
An unused good has multiplicity zero and contributes zero extra copies.
For integers $t,k\ge0$, a $(t,k)$-allocation satisfies
$\copies(X)\le t$ and $\extra_X(g)\le k$ for every $g\in G$.
Thus a cap of $k$ extra copies allows multiplicity at most $k+1$.

\begin{theorem}[Exact ordered-instance transfer]\label{thm:main}
Under the definitions above, let $A=(A_i)_{i\in N}$ be any allocation of
$\ranks$. Let $k\ge0$ be an integer and suppose
$\extra_A(a)\le k$ for every rank $a\in\ranks$.
There is an allocation $B=(B_i)_{i\in N}$ of $M$ such that
\begin{align}
 u_i(B_i)&\ge U_i(A_i) &&(i\in N),\label{eq:utility}\\
 \copies(B)&=\copies(A),\label{eq:copies}\\
 \extra_B(g)&\le k &&(g\in M),\label{eq:cap}\\
 \mult_B(g)&\ge1 &&(g\in M).\label{eq:cover}
\end{align}
The input may leave ranks unused. There is no restriction on which ranks
are copied or on their recipients.
\end{theorem}

The proof is given in Section~\ref{sec:proof}. The identity of copied goods
need not be preserved. Neither must the original bundle sizes $|A_i|$ be
preserved. In particular, completion can increase a bundle's size.

\begin{corollary}[Budgets and utility thresholds]\label{cor:threshold}
If $A$ is a $(t,k)$-allocation, the output in Theorem~\ref{thm:main} is a
$(t,k)$-allocation. If real numbers $\tau_i$ satisfy $U_i(A_i)\ge\tau_i$ for
all $i\in N$, the output also satisfies $u_i(B_i)\ge\tau_i$ for all agents.
\end{corollary}
\begin{proof}
Equation~\eqref{eq:copies} gives $\copies(B)=\copies(A)\le t$, and
\eqref{eq:cap} gives the per-good budget. The threshold assertion follows
from~\eqref{eq:utility}.
\end{proof}

\Needspace{12\baselineskip}
\paragraph{Why completion matters.}
Take two goods $a,b$ and two agents. Their weights are
\[
 (v_1(a),v_1(b))=(2,1),\qquad (v_2(a),v_2(b))=(1,2).
\]
Both ordered weight lists are $(2,1)$. If $A_1=A_2=\{1\}$, both ordered
utilities equal $2$ and the total extra-copy count is $1$.
Giving $a$ to agent~1 and $b$ to agent~2 achieves those utilities but uses
no extra copy. To preserve the count and cover both goods, the theorem may
instead give $B_1=\{a,b\}$ and $B_2=\{b\}$. Their utilities are $3$ and $2$,
and exactly one extra copy is used. This example also shows why the theorem
does not promise to preserve original bundle sizes.

\section{The rank-counting inequality}\label{sec:count}
The counting step converts rank information into a lower bound on eligible
agent--good incidences. It is the estimate needed for Hall's condition.

Suppose an ordered allocation $C$ covers every rank. Define its slot set by
\[
 \mathcal S=\{(i,a):i\in N,\ a\in C_i\},
 \qquad D=|\mathcal S|=\sum_{i\in N}|C_i|.
\]
An incidence $(i,g)\in N\times M$ is eligible for slot $(i,a)$ when
$r_i(g)\le a$. For $P\subseteq\mathcal S$, let $\nbr(P)$ be the set of all
incidences eligible for at least one slot in $P$. For $0\le s\le m$, put
\[
 c(s)=|\{(i,a)\in\mathcal S:a\le s\}|.
\]
This counts all slots in the prefix, including slots outside $P$.

\begin{lemma}\label{lem:count}
For every $P\subseteq\mathcal S$ and $I\subseteq N\times M$, let
$S=\{g\in M:\exists i,\ (i,g)\in I\}$ and $s=|S|$. Then
\begin{equation}\label{eq:rankcount}
 |P|\le |\nbr(P)\setminus I|+c(s).
\end{equation}
If every rank occurs between one and $q$ times in $C$, then also
\begin{equation}\label{eq:prefixbounds}
 c(s)\le qs,\qquad c(s)\le D-m+s.
\end{equation}
\end{lemma}
\begin{proof}
At most $c(s)$ slots in $P$ have rank at most $s$. Consider an agent $i$
with at least one selected slot of rank greater than $s$, and let $a_i$ be
its largest selected rank. Since $C_i$ is a set, the selected ranks greater
than $s$ are distinct members of $\{s+1,\ldots,a_i\}$. There are at most
$a_i-s$ of them.

The agent's top $a_i$ real goods form a set of size $a_i$, because $r_i$ is
a bijection. At most $s$ of those goods belong to $S$. Hence at least
$a_i-s$ of their incidences lie outside $I$. Each of these incidences is
eligible for the selected slot $(i,a_i)$, so it belongs to
$\nbr(P)\setminus I$. Incidences of different agents are distinct.
Summing over agents bounds the number of selected slots above $s$ by
$|\nbr(P)\setminus I|$. Adding the selected prefix slots proves
\eqref{eq:rankcount}.

There are $s$ prefix ranks, each occurring at most $q$ times, so $c(s)\le qs$.
Each of the last $m-s$ ranks occurs at least once. Thus at least $m-s$ of
the $D$ slots lie outside the prefix, which gives $c(s)\le D-(m-s)$.
\end{proof}

\section{A matching construction}\label{sec:proof}
The proof first covers every rank without changing the copy count. It then
uses a matching graph whose ports enforce the per-agent set condition and
whose primary vertices enforce coverage.

\begin{proof}[Proof of Theorem~\ref{thm:main}]
If $m=0$, all bundles are empty and the conclusions hold with $B_i=\varnothing$.
Assume $m\ge1$.

\paragraph{Complete the input.}
Fix an agent $i_0\in N$. Give each unused rank once to $i_0$, leaving all
other assignments unchanged, and denote the result by $C$.
Nonnegativity gives $U_i(C_i)\ge U_i(A_i)$ for every agent. An unused rank
changes multiplicity from zero to one, so its extra-copy count remains zero.
Consequently,
\[
 \copies(C)=\copies(A).
\]
Set $q=k+1$. Each rank now occurs between one and $q$ times. With $D$ as in
Section~\ref{sec:count}, double-counting incidences gives
\begin{equation}\label{eq:D}
 m\le D\le qm,\qquad \copies(C)=D-m.
\end{equation}

\paragraph{Build a balanced bipartite graph.}
The vertex classes are disjoint, even when their labels refer to the same
agent or good. The left and right sides have the following vertices.
\begin{center}
\begin{tabular}{@{}p{0.36\linewidth}rp{0.30\linewidth}r@{}}
\toprule
Left class & Count & Right class & Count\\
\midrule
Requested slots $(i,a)\in\mathcal S$ & $D$ & Right ports $(i,g)$ & $nm$\\
Left ports $(i,g)\in N\times M$ & $nm$ & Primary copies of goods & $m$\\
Dummy vertices & $qm-D$ & Secondary copies of goods & $(q-1)m$\\
\bottomrule
\end{tabular}
\end{center}
Each side has $nm+qm$ vertices. Every good has one primary and $q-1$
secondary copy vertices. The edges are exactly these:
\begin{enumerate}[nosep]
\item A slot $(i,a)$ is adjacent to right port $(i,g)$ if $r_i(g)\le a$.
\item Left port $(i,g)$ is adjacent to its identical right port and to all
$q$ copy vertices of good $g$.
\item Every dummy vertex is adjacent to every secondary copy vertex.
\end{enumerate}
In particular, slots and dummies have no primary-copy neighbors. Only the
identical left port and eligible slots can reach a given right port.

\paragraph{Verify Hall's condition.}
Take any subset of the left vertices. Write its slot part as $P$, its left
port part as the incidence set $I$, and its dummy part as $T$.
Set $S=\{g:\exists i,\ (i,g)\in I\}$ and $s=|S|$.
The right-port neighbors include $\nbr(P)\cup I$, of cardinality
\[
 |I|+|\nbr(P)\setminus I|.
\]
There are two cases.

If $T=\varnothing$, all $q$ copies of every good in $S$ are also neighbors.
They are distinct from right ports. The neighbor count is at least
\[
 |I|+|\nbr(P)\setminus I|+qs\ge |I|+|P|,
\]
by Lemma~\ref{lem:count} and $c(s)\le qs$.

If $T\ne\varnothing$, all $(q-1)m$ secondary copies are neighbors, as are
the $s$ primary copies indexed by $S$. The neighbor count is at least
\[
 |I|+|\nbr(P)\setminus I|+(q-1)m+s.
\]
Since $|T|\le qm-D$, the other bound in Lemma~\ref{lem:count} gives
\begin{align*}
 |P|+|T|
 &\le |\nbr(P)\setminus I|+(D-m+s)+(qm-D)\\
 &=|\nbr(P)\setminus I|+(q-1)m+s.
\end{align*}
Adding $|I|$ proves the required bound in this case too.

Thus every left-vertex subset has at least as many neighbors as vertices.
Hall's marriage theorem gives a matching saturating the left side.
The sides have equal finite cardinality, so the matching also saturates the
right side. It is a perfect matching.

\paragraph{Extract sets of goods.}
Each slot $(i,a)$ is matched to an eligible right port $(i,g)$.
Assign that good $g$ to agent $i$ and let $B_i$ be the resulting bundle.
Distinct slots use distinct right ports. In particular, two slots of the
same agent cannot assign the same good, so every $B_i$ is a set.

If a slot uses right port $(i,g)$, that right port cannot also be matched
through its identity edge. The left port $(i,g)$ must therefore be matched
to one of the $q$ copy vertices of $g$. Distinct assigned incidences use
distinct copy vertices. This proves $\mult_B(g)\le q$ for every good.

Every primary copy of $g$ has a matched left neighbor because the matching
is perfect. Such a neighbor must be a left port $(i,g)$.
This left port is occupied by the primary copy. Its right port cannot use
the identity edge and must therefore be matched from a slot.
That slot assigns $g$ to agent $i$. Hence $\mult_B(g)\ge1$ for all $g$.

All $D$ slots are filled, and no within-agent duplicate is removed in
forming the bundles. Thus $\sum_i|B_i|=D$. Coverage and double counting now
give
\[
 \copies(B)=\sum_{g\in M}(\mult_B(g)-1)=D-m
 =\copies(C)=\copies(A).
\]
Also $\extra_B(g)=\mult_B(g)-1\le q-1=k$.

Finally, eligibility gives $r_i(g)\le a$ for the good assigned to slot
$(i,a)$. Since the ordered values are nonincreasing,
\[
 v_i(g)=w_i(r_i(g))\ge w_i(a).
\]
Summing over agent $i$'s distinct slots yields
$u_i(B_i)\ge U_i(C_i)\ge U_i(A_i)$. This proves all four conclusions.
\end{proof}

\section{Algorithmic consequence and scope}
\begin{proposition}\label{prop:algorithm}
Given the rankings, a suitable allocation can be obtained by bipartite
matching on a graph with at most $2nm$ vertices on each side. For positive
$n,m$, a crude edge bound is $O(n^2m^2)$.
\end{proposition}
\begin{proof}
Set $\widehat q=\min\{k+1,n\}$. In a set-valued allocation, a rank occurs
at most $n$ times. The completed allocation therefore has multiplicities
at most $\widehat q$. Since $n\ge1$ and $k\ge0$, also $\widehat q\ge1$.
Repeat the construction with $q=\widehat q$.
Each side has $nm+\widehat q m\le2nm$ vertices. The slots contribute at
most $nm^2$ edges. Left ports contribute at most $nm(1+\widehat q)$ edges.
There are at most $nm$ dummies and at most $nm$ secondary copies, giving
at most $n^2m^2$ dummy edges. These counts imply the stated bound.
A standard bipartite maximum-matching algorithm finds a perfect matching,
and the proof describes how to read off the allocation.
\end{proof}

This is a polynomial number of combinatorial operations, even if $k$ is
large and binary encoded. With arbitrary real weights and supplied
rankings, no bit-complexity model for real-number comparisons is assumed.
No executable matching algorithm or runtime bound is asserted to have been
formalized in Lean.

Nonnegativity is used when unused ranks are assigned during completion.
The matching and rankwise utility argument for an already covered input
only requires real weights ordered by the supplied bijections.
The construction preserves $|C_i|$ for each completed bundle, but it need
not preserve $|A_i|$. It can also change which named goods are copied.
The theorem transfers achieved utility thresholds. It does not by itself
prove that one extra copy suffices for exact maximin-share fairness, the
separate fair-division question discussed in the source context.

\Needspace{12\baselineskip}
\section{Formal statement map and reproducibility}
The accompanying source targets Lean~4.19.0 and mathlib revision
\begin{center}
\texttt{\small c44e0c8ee63ca166450922a373c7409c5d26b00b}.
\end{center}
The three final declarations, in namespace
\texttt{OrderedReductionAllocation}, are
\begin{verbatim}
ordered_allocation_transfer_exact
ordered_allocation_transfer
ordered_threshold_transfer
\end{verbatim}
They correspond to Theorem~\ref{thm:main}, its budget consequence, and the
utility-threshold consequence. The definitions of multiplicity, extra
copies, total copies, and utility are in namespace
\texttt{OrderedReductionCompletion}. Natural-number subtraction there is
truncated at zero, so its extra-copy definition agrees with the maximum
used above. Lean uses zero-based \texttt{Fin m} ranks; this paper uses
one-based ranks. The eligibility order and all counts agree under that
change of labels.

The formal source includes the prefix estimates, the rank-counting
inequality, the Hall gadget and assignment extraction, completion, and the
allocation-level utility and exact-copy statements. A declaration map and reproduction instructions accompany this manuscript.
A fresh Lean 4.19.0 build succeeded for the import-only source variant,
which replaces the umbrella import by two explicit imports and leaves the
proof body unchanged. A separate audit with \texttt{-{}-trust=0} checked the
three final declarations and an expanded statement of the exact transfer.
All four reports contain only the standard foundational axioms
\texttt{propext}, \texttt{Classical.choice}, and \texttt{Quot.sound}.
A formal proof checks the stated proposition, not historical novelty or
external expert validation.

\bibliographystyle{plain}
\bibliography{references}
\end{document}
